\documentclass[conference]{IEEEtran}
\IEEEoverridecommandlockouts

\usepackage{cite}
\usepackage{amsmath,amssymb,amsfonts}
\usepackage{graphicx}
\usepackage{textcomp}
\usepackage{xcolor}
\usepackage{booktabs}
\usepackage{url}
\usepackage{multirow}
\usepackage{tikz}
\usetikzlibrary{arrows.meta, positioning}
\def\BibTeX{{\rm B\kern-.05em{\sc i\kern-.025em b}\kern-.08em
    T\kern-.1667em\lower.7ex\hbox{E}\kern-.125emX}}
\begin{document}

\title{Beyond Access: Explainable Machine Learning Reveals a Connected but Unengaged Population in Pakistan's Digital Financial System}

\author{
\IEEEauthorblockN{Abdur-Rehman Ashar}
\IEEEauthorblockA{
Beaconhouse School System PTC \\
Gujranwala, Pakistan \\
abdurrehmanashar@gmail.com
}
\and
\IEEEauthorblockN{Umema Ashar}
\IEEEauthorblockA{
KDD Lab, FAST NUCES \\
Islamabad, Pakistan \\
umema2004@gmail.com
}
}

\maketitle
\begin{abstract}
Financial inclusion policy in Pakistan, as in much of the developing world,
is typically measured by formal account ownership --- but ownership does
not tell us whether people actually use digital financial tools once they
have access. Using individual-level data from the Global Findex Database
2025 (Pakistan, $n=1{,}000$), this study finds that \textbf{44.3\% of the
Findex Pakistan analytical sample} belongs to a distinct segment we call ``Connected but
Unengaged'': near-universal mobile phone ownership (98.0\%) and majority
internet use (62.3\%), yet two-thirds (66.8\%) remain formally excluded
from the financial system --- nearly ten times the directly-measured
dormant-account rate (6.6\%). This segment is not an artifact of a single
method: it emerges independently from unsupervised clustering, from the
population profile of the classifier's prediction errors, and from a
feature-group ablation study showing behavioural variables, not access
variables, drive adoption in this context. The convergence of three
unrelated analytical approaches on the same population boundary is the
paper's central evidence that, for a large share of Pakistanis, this suggests
that connectivity alone does not explain non-adoption and points toward downstream
barriers -- such as trust, product design, or institutional friction -- that are not captured in survey data. We support this finding with an
explainable ML pipeline (Random Forest, XGBoost, SHAP; macro F1 = 0.494,
validated via repeated cross-validation) and a bootstrap-validated cluster
structure (Adjusted Rand Index = 0.904). We also find that gender, not
urban/rural status, is the strongest structural predictor of segment
membership (Cramer's $V = 0.807$), and that naive subgroup accuracy
comparisons are misleading once gender's entanglement with outcome is
accounted for --- a methodological caveat likely to generalize beyond this
setting. All findings are associational; the cross-sectional nature of the
data does not support causal claims about why this segment remains
unengaged.
\end{abstract}

\begin{IEEEkeywords}
fintech adoption, financial inclusion, explainable AI, SHAP, k-means
clustering, model calibration, algorithmic fairness, Pakistan, Global Findex
\end{IEEEkeywords}
 
\section{Introduction}
\label{sec:introduction}

A dormant-account rate of 6.6\% captures inactivity among people who already
hold an account, but does not capture the potentially larger population who
have digital access yet never enter the formal financial system. This paper shows that number
is almost certainly the wrong denominator. Using individual-level Findex
2025 data, we find that when access is measured properly --- mobile phone
ownership and internet use, rather than account status alone --- a much
larger population (44.3\% of the analytical sample) has the technical capacity to
engage digitally but does not. This gap is invisible to the standard
dormant-account metric, which only counts people who already hold an
account; it misses everyone who has connectivity but never obtained one in
the first place. We surface this segment independently through three unrelated analytical
methods and show the pattern holds under model-error inspection,
feature-group ablation, and bootstrap validation. Establishing that this
population exists and is measurable is the contribution of this paper;
explaining why its members remain unengaged despite access is left as
the natural next question.

This study applies explainable machine learning and unsupervised
segmentation to individual-level survey data to identify who is excluded
from Pakistan's digital financial system and, more specifically, where
formal access fails to translate into actual digital usage. Using Random
Forest and XGBoost classifiers with SHAP explainability alongside $k$-means
clustering, the analysis tests three hypotheses concerning the access-usage
gap, the urban-rural adoption divide, and the relative importance of
demographic versus access-based predictors of adoption. Beyond the core
predictive and segmentation results, we subject every major finding to an
explicit robustness check: hyperparameter-tuned models are evaluated under
repeated cross-validation rather than a single train-test split; the
SHAP-based feature-importance ranking is cross-checked against a feature-group ablation
study using an entirely different methodology; and the cluster structure is
validated via bootstrap resampling and an algorithm-independent comparison
against agglomerative clustering. To our knowledge, we found no prior South 
Asian study combining these analytical components with the same robustness and validation framework.

\subsection{Novelty Statement}
\label{subsec:novelty_statement}

This work makes four specific contributions beyond prior literature.

First, it identifies and triangulates --- via clustering, model-error
analysis, and ablation --- a large access-rich, adoption-poor population
segment that is substantially larger than the directly-measured
dormant-account rate would suggest, demonstrating the value of combining
multiple analytical lenses rather than relying on a single metric.

Second, it provides an explicit robustness and calibration audit rarely
reported in this literature, showing that headline accuracy figures can
mask systematic underconfidence for the most policy-relevant class
(Excluded) and that naive subgroup accuracy comparisons can be misleading
when a demographic variable is strongly entangled with outcome.

Third, it offers a direct, quantitative test of whether the access-dominant predictor ranking reported for Sub-Saharan Africa~\cite{saidy2025} generalizes
to Pakistan, finding that it does not --- behavioural and demographic
variables outrank raw access variables in this context, a result confirmed
independently through SHAP, permutation importance, and ablation.

Fourth, it extends the dual supervised-unsupervised
explainable ML methodology of Chowdhury et al.~\cite{chowdhury2026} to
South Asian survey data, using the newer 2024-fieldwork Findex round rather
than the 2021 round used in comparable Sub-Saharan African
studies~\cite{saidy2025}.

\subsection{Research Questions}
\label{subsec:research_questions}

The analysis is organized around four research questions, which the
hypotheses and analyses in Sections~\ref{sec:data_and_methods}--\ref{sec:results} address individually and, in
combination, jointly.

\begin{itemize}
\item[\textbf{RQ1.}] To what extent does formal financial access translate
into digital financial usage among Pakistani adults, and does this
translation vary by urban/rural context?
\item[\textbf{RQ2.}] Can unsupervised learning, applied independently of
any adoption label, identify distinct behavioural/access population
profiles --- and if so, do these profiles reveal structure not visible in
directly-measured adoption statistics (e.g., the dormant-account rate)?
\item[\textbf{RQ3.}] Which demographic, access, behavioural, and
resilience variables are most strongly associated with differences in
Adoption Tier, and does the resulting ranking generalize across
explanation methods (SHAP, Gini, permutation) and across national
contexts (cf.\ Sub-Saharan Africa~\cite{saidy2025})?
\item[\textbf{RQ4.}] Are the population segments identified by clustering
robust --- to resampling, to the choice of clustering algorithm, and to
corroboration through independent diagnostic routes (model error
analysis, feature-group ablation)?
\end{itemize}

Table~\ref{tab:rq_map} maps each research question to the hypotheses
tested in Section~\ref{sec:results} and the methods used to address it. RQ1 and RQ2 are
addressed jointly: RQ1 through the directly-measured access-usage gap
(H1) and the urban/rural comparison (H2), and RQ2 through the independent
unsupervised segmentation that corroborates and extends the H1 finding
(Section~\ref{subsec:h1_access_usage_gap_rq1}). RQ3 corresponds to H3, the supervised modelling and
explainability analysis. RQ4 is not tested via a numbered hypothesis;
instead, it is addressed directly through the cluster validation suite
(Section~\ref{subsec:cluster_validation_results}) --- bootstrap resampling, the Ward agglomerative
cross-check, and the $k=2$ vs.\ $k=3$ comparison --- and through the
convergence of clustering, error analysis, and ablation summarized in
Table~\ref{tab:convergence}.

\begin{table}[h]
\caption{Mapping of Research Questions to Hypotheses and Methods}
\label{tab:rq_map}
\centering
\renewcommand{\arraystretch}{1.3}
\begin{tabular}{p{0.7cm}p{3.4cm}p{3.4cm}}
\toprule
\textbf{RQ} & \textbf{Addressed by} & \textbf{Section(s)} \\
\midrule
RQ1 & H1 (access-usage gap), H2 (urban/rural divide) & \S\ref{subsec:h1_access_usage_gap_rq1}, \S\ref{subsec:h2_urban_versus_rural_adoption_rq1} \\
\addlinespace
RQ2 & Unsupervised $k$-means segmentation, independent of Adoption Tier &
\S\ref{subsec:unsupervised_segmentation}, \S\ref{subsec:h1_access_usage_gap_rq1} \\
\addlinespace
RQ3 & H3 (supervised modelling); SHAP, Gini, permutation importance &
\S\ref{subsec:supervised_modelling}--\S\ref{subsec:feature_importance_consistency_methods}, \S\ref{subsec:h3_predictors_of_adoption_rq3} \\
\addlinespace
RQ4 & Bootstrap ARI, Ward cross-check, $k=2$ vs.\ $k=3$ comparison,
three-method convergence (Table~\ref{tab:convergence}) & \S\ref{subsec:cluster_validation_results}, \S\ref{sec:discussion} \\
\bottomrule
\end{tabular}
\end{table}

\section{Related Work}
\label{sec:related_work}
 
\subsection{Measuring Financial Inclusion}
\label{subsec:measuring_financial_inclusion}
The Global Findex Database, now in its fifth edition, remains the primary demand-side data source for cross-country financial inclusion research \cite{demirguckunt2022}. Early editions of Findex research documented steady global growth in account ownership but repeatedly noted that account ownership alone is an incomplete proxy for financial inclusion, since it does not capture whether an account is actively used. Suri and Jack \cite{surijack2016} provided some of the most influential causal evidence on this front, showing via a long-run panel study in Kenya that access to mobile money -- specifically, active usage rather than mere availability -- reduced extreme poverty and produced disproportionate gains for female-headed households through occupational shifts into business and retail activity. Their finding that gendered impacts of digital financial services can be large and structurally significant directly motivates our attention to gender as a candidate predictors of modelled fintech adoption, rather than treating it as a routine demographic control variable.
 
\subsection{Machine Learning Approaches to Fintech Adoption}
\label{subsec:ml_approaches_fintech_adoption}
Three recent studies directly inform the design of this project. Chowdhury et al. \cite{chowdhury2026} applied Random Forest, XGBoost, and k-means clustering to financial literacy and behaviour data from 1{,}067 adults in Bangladesh, validating their segmentation with silhouette, Davies-Bouldin, and Calinski-Harabasz indices. They identified three behavioural clusters (34\%, 41\%, and 25\% of the sample) and reported a weak correlation between financial knowledge and financial behaviour, alongside a counterintuitive 7.3 percentage-point rural literacy advantage over urban respondents, with formal education emerging as a comparatively weak predictor relative to institutional trust.
 
Emuron, Kapim Kenfack, and Msimanga \cite{emuron2025} took a more applied approach, predicting fintech adoption in South Africa from Findex, mobile pricing, and central bank data using regression-based methods on remittances, borrowing, age, and digital usage variables. Their work established a directly comparable variable set to the one constructed here from Findex Pakistan, though without a segmentation component or an explicit robustness audit.
 
Saidy and Hassan \cite{saidy2025} modelled fintech adoption drivers across 36 Sub-Saharan African countries using the 2021 Global Findex round, combining Logit and Probit regression with Random Forest, Gradient Boosting, and XGBoost. Their central finding was that formal account ownership and technology access -- specifically mobile phone ownership and internet connectivity -- were the strongest predictors of adoption, a finding we directly test for generalizability in Section~\ref{sec:discussion}.
 
\subsection{Positioning}
\label{subsec:positioning}
This project combines Chowdhury et al.'s dual supervised-unsupervised methodology with the adoption-prediction framing of Emuron et al. and Saidy and Hassan, applying both to Pakistan using the newer 2024-fieldwork Findex round \cite{worldbank2025} -- a combination not yet applied to South Asian data. It further extends all three prior works methodologically by adding hyperparameter-tuned robustness estimates, SHAP dependence analysis, calibration diagnostics, cross-method feature importance validation, model error analysis linked back to the segmentation results, an ablation study, and bootstrap-validated cluster stability -- none of which are reported in the three anchor papers.
 
\section{Data and Methods}
\label{sec:data_and_methods}
 
\subsection{Data}
\label{subsec:data}
The analysis uses the Global Findex Database 2025, Pakistan microdata (Reference: PAK\_2024\_FINDEX\_v02\_M) \cite{worldbank2025}, an individual-level survey of 1{,}000 respondents conducted by the World Bank between May and December 2024. The dataset includes 183 variables spanning demographics, account ownership, digital and informal financial behaviour, and financial resilience. Findex microdata is collected under informed consent protocols administered by the World Bank and its survey partners and is distributed in de-identified form; no personally identifying information is available to or used by this analysis (see Section~\ref{subsec:data_privacy} for further discussion). \textbf{All
analyses in this paper use unweighted respondent counts; reported
proportions describe the sample as collected and are not adjusted for
Findex's survey design weights (Section~VII). This choice reflects the study's focus on predictive modelling and behavioural segmentation within the observed analytical sample rather than estimation of population-level prevalence; consequently, reported proportions are interpreted as sample descriptors.}

 
\subsection{Target Construction}
\label{subsec:target_construction}
A four-level ordinal Adoption Tier was constructed from account ownership and digital usage indicators: Excluded (no formal account), Basic (account held, no digital usage), Digitally Engaged (uses digital payments and/or a digitally enabled account), and Advanced (digital usage combined with formal saving or borrowing behaviour). The resulting distribution was 68.2\% Excluded, 2.1\% Basic, 10.1\% Digitally Engaged, and 19.6\% Advanced, shown in Fig.~\ref{fig:tier_overview}.
 
\subsection{Feature Selection and Leakage Control}
\label{subsec:feature_selection_leakage_control}

Twenty variables across four construct groups (demographics, access,
behaviour, resilience) were initially selected from the Findex codebook as
candidate predictors.

\textbf{Leakage audit.} An initial model run retaining all twenty
candidate variables --- including \texttt{account}, \texttt{dig\_account},
\texttt{anydigpayment}, \texttt{saved}, and \texttt{fin22a} --- produced
spurious near-perfect classification accuracy. Inspection traced this to
target leakage: the Adoption Tier label was constructed directly from
these five variables (Section~\ref{subsec:target_construction}), and \texttt{account\_fin} and
\texttt{account\_mob} are near-deterministic components of
\texttt{account} itself. In effect, the initial specification allowed the
model to partially reconstruct its own target from near-identical inputs.
\textbf{We report this failed initial specification deliberately, to
demonstrate that leakage was detected and corrected before final
evaluation}, rather than omitting it as if the leakage-safe specification
had been the only one attempted.

All seven implicated variables were removed, leaving fifteen leakage-safe
predictors organized into four construct groups (Table~\ref{tab:feature_groups}). This
fifteen-variable set is the one used throughout the remainder of the
paper --- for supervised modelling (Section~\ref{subsec:supervised_modelling}), SHAP and permutation
importance (Sections~\ref{subsec:explainability}, \ref{subsec:feature_importance_consistency_methods}), ablation (Section~\ref{subsec:ablation_study_methods}), and
unsupervised clustering (Section~\ref{subsec:unsupervised_segmentation}). Fig.~\ref{fig:leakage_pipeline}
summarizes the audit sequence.

\begin{figure}[h]
\centering
\fbox{\begin{minipage}{0.9\linewidth}
\centering
\textbf{Initial candidate feature set} (20 variables) \\
$\big\downarrow$ \\
\textit{Leakage audit --- near-perfect accuracy flagged as spurious} \\
$\big\downarrow$ \\
\textbf{7 target-derived variables removed} \\
\footnotesize (\texttt{account}, \texttt{dig\_account},
\texttt{anydigpayment}, \texttt{saved}, \texttt{fin22a},
\texttt{account\_fin}, \texttt{account\_mob}) \\
\normalsize $\big\downarrow$ \\
\textbf{15 leakage-safe predictors} (Table~\ref{tab:feature_groups}) \\
$\big\downarrow$ \\
\textbf{Final modelling} \\
\footnotesize (supervised classification, SHAP/permutation/ablation,
$k$-means clustering)
\end{minipage}}
\caption{Leakage audit and correction sequence. The initial specification
included target-derived variables and produced spurious near-perfect
accuracy; this was identified, corrected, and is reported here for
methodological transparency.}
\label{fig:leakage_pipeline}
\end{figure}

\begin{table}[h]
\caption{Feature Groups Used in Modelling and Ablation}
\label{tab:feature_groups}
\centering
\renewcommand{\arraystretch}{1.3}
\begin{tabular}{@{}p{2.1cm}p{5.7cm}@{}}
\toprule
\textbf{Group} & \textbf{Variables} \\
\midrule
Demographics & \texttt{female}, \texttt{age}, \texttt{educ},
\texttt{inc\_q}, \texttt{in\_workforce}, \texttt{is\_rural} \\
\addlinespace
Access & \texttt{con1}, \texttt{internet\_use} \\
\addlinespace
Behaviour & \texttt{merchantpay\_dig}, \texttt{fin17b},
\texttt{fin17c}, \texttt{borrowed}, \texttt{fin22a\_1} \\
\addlinespace
Resilience & \texttt{fin24a}, \texttt{fin24b} \\
\bottomrule
\end{tabular}
\end{table}
 
\subsection{Modelling Objective}
\label{subsec:modelling_objective}

Before presenting the supervised pipeline, we state its intended role
explicitly. \textbf{The objective of supervised modelling in this study
was not to develop a deployment-ready individual-level classifier, but to
identify and interrogate the behavioural and demographic correlates of
adoption} through explainability tooling (SHAP, permutation importance)
and diagnostic analysis (calibration, error profiling, ablation). Given a
four-level target with substantial class imbalance (68.2\% Excluded,
2.1\% Basic) and $n=1{,}000$, a macro F1 in the 0.45--0.50 range is the
realistic ceiling for this feature set and sample size, not a shortfall to
be corrected by further tuning. We report this ceiling candidly and treat
the fitted model as an analytical instrument --- a lens for extracting
feature-outcome relationships and cross-validating the unsupervised
segmentation in Section~\ref{subsec:h1_access_usage_gap_rq1} --- rather than as an endpoint whose accuracy
is itself the contribution.

\subsection{Supervised Modelling}
\label{subsec:supervised_modelling}

Logistic Regression, Decision Tree, Random Forest, and XGBoost classifiers
were trained as candidate models for extracting adoption correlates, not
as competing deployment candidates. All were trained on a stratified
75/25 train-test split, with class weighting applied to address the
imbalanced target distribution. The best-performing default configuration
(Random Forest) was then re-tuned via \texttt{RandomizedSearchCV} (25
iterations, 5-fold stratified cross-validation, scored on macro F1) over a
search space covering the number of estimators, maximum depth, minimum
samples per leaf, and maximum features considered per split; the full
search space is reported in Appendix~\ref{sec:appendix_hyperparameter_search_space}. An equivalent search was run for
XGBoost for comparison. Tuning aimed to obtain a stable, well-characterized
model for explainability analysis, not to maximize headline accuracy.

\subsection{Model Robustness}
\label{subsec:model_robustness_methods}
To assess whether performance estimates were sensitive to the particular train-test split drawn, the tuned Random Forest was additionally evaluated under repeated stratified 5-fold cross-validation with 10 repeats (50 total folds), yielding a distribution of macro F1 scores rather than a single point estimate.
 
\subsection{Explainability}
\label{subsec:explainability}

The tuned Random Forest was explained using SHAP TreeExplainer values.
Because the target is multiclass, per-class SHAP arrays were averaged
(mean absolute value) to produce an overall importance ranking, and SHAP
dependence plots were generated for the top four ranked features against
the Advanced-tier class specifically, as this class is the most directly
policy-relevant ``successful adoption'' outcome.

\textbf{Terminology.} Throughout this paper we refer to SHAP, Gini, and
permutation rankings as identifying \emph{predictors} or \emph{influential
features} --- terms describing a variable's contribution to model output
--- rather than \emph{predictors} or \emph{causes} of adoption. This
distinction matters because Findex is a cross-sectional survey (Section
III-A) and no causal identification strategy (e.g., instrumental
variables, panel data, or experimental variation) is employed here. A
feature ranked highly by SHAP is one the model relies on to make
predictions; this is informative about association and about which
factors are worth investigating further, but it does not establish that
changing that feature would change adoption behaviour.
 
\subsection{Calibration Analysis}
\label{subsec:calibration_analysis_methods}
Predicted class probabilities were evaluated for calibration using one-vs-rest reliability diagrams (5 quantile-based bins per class) and the Brier score, to assess whether predicted probabilities correspond to observed outcome frequencies -- a property distinct from, and not guaranteed by, classification accuracy.
 
\subsection{Feature Importance Consistency}
\label{subsec:feature_importance_consistency_methods}
To assess whether the SHAP-based feature-importance ranking was an artifact of a single explanation method, we additionally computed Gini importance (native to the Random Forest) and permutation importance (20 repeats, scored on macro F1 degradation on the held-out test set), and compared all three rankings via Spearman rank correlation.
 
\subsection{Error Analysis}
\label{subsec:error_analysis_methods}
Misclassified test-set observations were examined for (i) the dominant confusion pair, (ii) subgroup accuracy by gender and urban/rural status, and (iii) the mean feature profile of misclassified versus correctly classified observations within the dominant confusion pair, to determine whether model errors carried substantive meaning rather than representing unstructured noise.
 
\subsection{Ablation Study}
\label{subsec:ablation_study_methods}
To independently validate the SHAP-based ranking of construct groups, we conducted a feature-group ablation study, retraining the tuned Random Forest configuration on (i) the full feature set, (ii) each construct group individually held out, and (iii) each construct group used in isolation, measuring macro F1 in each configuration.
 
\subsection{Unsupervised Segmentation}
\label{subsec:unsupervised_segmentation}

$K$-means clustering was applied to the same leakage-safe feature set
(standardized via $z$-score) used in supervised modelling --- the fifteen
demographic, access, behaviour, and resilience variables in Table~\ref{tab:feature_groups} --- for
$k=2$ to $5$, with $k$ selected by silhouette score and Davies-Bouldin and
Calinski-Harabasz indices computed for robustness (Fig.~\ref{fig:k_selection}).

\textbf{The Adoption Tier target variable was not included as a clustering
input under any configuration.} Clustering was performed entirely on
demographic, access, behavioural, and resilience predictors; Adoption Tier
was used exclusively \emph{after} cluster assignment, to profile and
interpret the resulting groups (Table~\ref{tab:clusters}). This separation is deliberate:
because Adoption Tier is itself constructed from account and digital-usage
indicators (Section III-B), including it as a clustering feature would
risk the clusters trivially reproducing the target rather than revealing
independent structure. Clusters were profiled on demographic and access
characteristics and visualized via a two-component PCA projection
(Fig.~\ref{fig:pca_scatter}); the resulting cluster labels (e.g., ``Connected but Unengaged'')
are therefore \emph{descriptive summaries applied post hoc}, not
categories the clustering algorithm was told to find.
 
\paragraph{Defining ``Connected but Unengaged.''}
We use this label to describe a cluster whose members are, on average,
\emph{connected} by conventional access metrics (mobile phone ownership,
internet use) but \emph{unengaged} by the independently-constructed
Adoption Tier (i.e., predominantly in the Excluded tier despite that
access). The label names a statistical regularity observed after
clustering, not a variable the algorithm optimized for.

\subsection{Cluster Validation}
\label{subsec:cluster_validation_methods}
Three checks were used to validate the k-means solution beyond the silhouette 
score alone. First, bootstrap stability was assessed by refitting k-means 
on 100 bootstrap resamples and computing the Adjusted Rand Index (ARI) 
between each resample's cluster assignment (applied to the full original 
sample) and the original clustering. Second, an algorithm-independent 
cross-check compared the k-means solution against Ward agglomerative 
clustering at the same $k$. Third, because the silhouette score for 
$k=3$ was only marginally higher than for $k=2$ (0.1675), 
we conducted an explicit interpretability comparison between the two 
solutions, examining what population-level distinction the additional 
cluster does or does not preserve. We therefore interpret the three-cluster 
solution as a substantively useful segmentation of the observed sample rather 
than evidence that three sharply separated natural populations exist.
 
\subsection{Statistical Validation}
\label{subsec:statistical_validation}
Chi-square tests of independence with Cramer's V effect sizes were used to test associations between cluster membership, Adoption Tier, gender, and urban/rural status. Wilson score 95\% confidence intervals were computed for all headline proportions, given the small size of some subgroups (e.g. $n=21$ for the Basic tier).
 
\begin{figure}[h]
\centering
\includegraphics[width=\linewidth]{adoption_tier_overview.png}
\caption{Adoption Tier distribution, H1 access-usage gap, and Adoption Tier by urban/rural and income quintile.}
\label{fig:tier_overview}
\end{figure}

\section{Results}
\label{sec:results}
 
\subsection{H1 -- The Access-Usage Gap (RQ1)}
\label{subsec:h1_access_usage_gap_rq1}
Among the 318 respondents holding a formal account, 6.6\% (95\% CI: 4.4\%--9.9\%) showed no digital financial usage -- a direct measure of the dormant-account phenomenon described by Chowdhury et al. \cite{chowdhury2026} as a knowledge-behaviour gap in the Bangladesh context.
 
A substantially larger and more striking pattern emerged from the
independent, unsupervised cluster analysis. $K$-means ($k=3$), fitted
\emph{without access to the Adoption Tier label} (see Section~\ref{subsec:unsupervised_segmentation}),
identified a cluster --- labelled ``Connected but Unengaged'' after the
fact based on its profile --- comprising 44.3\% of the entire sample
($n=443$). This cluster shows near-universal mobile phone ownership
(98.0\%) and majority internet use (62.3\%), yet 66.8\% (95\% CI:
62.3\%--71.0\%) remain formally Excluded from the financial system
(Table~\ref{tab:clusters}). Because this pattern emerged from clustering on access and
demographic variables with no reference to the Adoption Tier during
fitting, it constitutes independent confirmation that access is not the
binding constraint for a large share of the population.

\subsection{H2 -- Urban versus Rural Adoption (RQ1)}
\label{subsec:h2_urban_versus_rural_adoption_rq1}
Adoption Tier was significantly associated with urban/rural status ($\chi^2 = 21.742$, $p = 0.0001$, $\mathit{df} = 3$, Cramer's V $= 0.147$), and cluster membership showed a parallel association ($\chi^2 = 19.823$, $p < 0.0001$, $\mathit{df} = 2$, Cramer's V $= 0.141$). The direction of the effect favoured urban respondents: rural share declined monotonically across clusters of increasing adoption, from 51.0\% in the least-adopted cluster to 46.0\% and then 28.6\% in progressively more digitally engaged clusters. This is consistent with the conventional urban-adoption advantage reported broadly in the financial inclusion literature, and contrasts with Chowdhury et al.'s reversed finding of a rural literacy advantage in Bangladesh \cite{chowdhury2026}.
 
\subsection{H3 -- Predictors of Adoption (RQ3)}
\label{subsec:h3_predictors_of_adoption_rq3}
 
\begin{table}[h]
\caption{Supervised Model Comparison}
\label{tab:models}
\centering
\begin{tabular}{@{}lccc@{}}
\toprule
Model & Accuracy & F1 (macro) & F1 (weighted) \\
\midrule
XGBoost & 0.800 & 0.481 & 0.781 \\
Random Forest (default) & 0.708 & 0.486 & 0.737 \\
Random Forest (tuned) & -- & \textbf{0.494} & -- \\
Decision Tree & 0.488 & 0.424 & 0.554 \\
Logistic Regression & 0.420 & 0.392 & 0.473 \\
\bottomrule
\end{tabular}
\end{table}
 
\begin{figure}[h]
\centering
% TODO: add confusion_matrix_best_model.png here
\includegraphics[width=\linewidth]{confusion_matrix_best_model.png}
\caption{Confusion matrix for the best-performing model (Random Forest).}
\label{fig:confusion_matrix}
\end{figure}
 
SHAP analysis of the tuned Random Forest (Fig.~\ref{fig:shap_summary}) ranked predictors as follows: saving via mobile money (\texttt{fin17b}), age, workforce participation, informal saving (\texttt{fin17c}), income quintile, mobile phone ownership, emergency fund coverage time, gender, and internet use. Mobile phone ownership and internet use -- the two strongest predictors identified by Saidy and Hassan across Sub-Saharan Africa \cite{saidy2025} -- ranked in the middle of the importance distribution rather than at the top.

\begin{figure}[h]
\centering
% TODO: add feature_importance.png here
\includegraphics[width=\linewidth]{feature_importance.png}
\caption{Random Forest feature importance (Gini-based).}
\label{fig:feature_importance}
\end{figure}
 
\begin{figure}[h]
\centering
\includegraphics[width=\linewidth]{shap_summary_tuned.png}
\caption{Mean absolute SHAP value per feature, averaged across all four Adoption Tier classes.}
\label{fig:shap_summary}
\end{figure}
 
\subsubsection{SHAP Dependence Plots}
\label{subsubsec:shap_dependence_plots}
Fig.~\ref{fig:shap_dependence} shows dependence plots for the four top-ranked features against the Advanced-tier class. \texttt{fin17b} and \texttt{fin17c} show clean binary step functions, as expected for binary indicators. Notably, \texttt{age} shows a nonlinear relationship: SHAP contribution rises through the twenties and thirties, peaks around ages 40--50, and declines thereafter -- indicating that the model's attribution of age to predicted adoption is not simply monotonic, with the largest positive contribution concentrated among economically active, middle-aged respondents.

\begin{figure}[h]
\centering
\includegraphics[width=\linewidth]{shap_dependence_plots.png}
\caption{SHAP dependence plots for the top four features (fin17b, age, in\_workforce, fin17c) against the Advanced-tier class, colored by the strongest interacting feature.}
\label{fig:shap_dependence}
\end{figure}
 
\subsubsection{Calibration Analysis}
\label{subsubsec:calibration_analysis_results}
Fig.~\ref{fig:calibration} shows per-class reliability diagrams. The Excluded class is poorly calibrated and systematically underconfident (Brier $=0.168$): when the model assigns roughly 40\% probability to Excluded, the true observed frequency is closer to 65--70\%. The Advanced class is comparatively well-calibrated (Brier $=0.042$), closely tracking the diagonal. The Basic and Digitally Engaged classes show jagged, unreliable calibration curves, consistent with their small effective sample sizes.
 
\begin{figure}[h]
\centering
\includegraphics[width=\linewidth]{calibration_reliability_diagrams.png}
\caption{Per-class reliability diagrams (one-vs-rest) with Brier scores.}
\label{fig:calibration}
\end{figure}

\subsubsection{Feature Importance Consistency}
\label{subsubsec:feature_importance_consistency_results}
Table~\ref{tab:importance_consistency} and Fig.~\ref{fig:importance_consistency} compare Gini, permutation, and SHAP importance rankings. Gini and SHAP rankings agree strongly (Spearman $\rho = 0.914$, $p < 0.0001$); permutation importance agrees more weakly with SHAP ($\rho = 0.504$, $p = 0.056$, not significant at $\alpha=0.05$) and is negative for several features (\texttt{con1}, \texttt{female}, \texttt{internet\_use}, \texttt{fin24a}, \texttt{educ}, \texttt{is\_rural}), indicating that shuffling these features does not reliably degrade -- and can marginally improve -- macro F1 on the test set. \texttt{fin17b} is the only feature ranked at or near the top by all three methods, making it the single most robust predictive feature identified in this analysis, subject to the routing-pattern caveat discussed in Section~\ref{sec:limitations}.

\begin{table}[h]
\caption{Top Features by Cross-Method Importance (abbreviated; full table in Appendix~\ref{sec:appendix_additional_diagnostic_tables})}
\label{tab:importance_consistency}
\centering
\begin{tabular}{@{}lccc@{}}
\toprule
Feature & Gini rank & Permutation rank & SHAP rank \\
\midrule
fin17b & 2 & 1 & 1 \\
age & 1 & 7 & 2 \\
in\_workforce & 5 & 3 & 3 \\
fin17c & 6 & 2 & 4 \\
inc\_q & 3 & 6 & 5 \\
\bottomrule
\end{tabular}
\end{table}

Model comparison results (Table~\ref{tab:models}) show XGBoost achieved the highest
single-split accuracy (0.800) but Random Forest achieved the highest macro
F1 (0.486 default; 0.494 tuned), the metric prioritized here given
substantial class imbalance in the target. Consistent with the modelling
objective stated in Section~\ref{subsec:modelling_objective}, this level of performance is modest
in absolute terms and should be read as adequate for extracting directional
feature-outcome relationships, not as evidence of a strong individual-level
predictive model. All models struggled with the Basic tier ($n=21$),
attributable to sample size rather than model deficiency (Fig.~\ref{fig:confusion_matrix}); given
this small subgroup, per-class metrics for the Basic tier should be
treated as indicative only and are not a meaningful test of model quality.
 
\begin{figure}[h]
\centering
% TODO: add feature_importance_consistency.png here
\includegraphics[width=\linewidth]{feature_importance_consistency.png}
\caption{Normalized feature importance across Gini, permutation, and SHAP methods.}
\label{fig:importance_consistency}
\end{figure}
 
\subsubsection{Error Analysis}
\label{subsubsec:error_analysis_results}
The dominant confusion pair was Excluded predicted as Digitally Engaged (21 cases), with the reverse error (Digitally Engaged predicted as Excluded, 16 cases) also common -- consistent with this boundary being the hardest to separate given the available features.
 
Subgroup accuracy was 0.852 for women versus 0.689 for men (ratio 0.808), and 0.780 for rural versus 0.752 for urban respondents (ratio 0.964). We caution against interpreting the gender gap as evidence the model treats women ``more fairly'': because gender is very strongly associated with cluster membership and outcome in this population (Cramer's V $=0.807$; see Section~\ref{subsec:cluster_profiles_gender_association}), the majority of women fall into the large, easy-to-predict Excluded class, mechanically inflating aggregate accuracy for that subgroup. We regard this as an important methodological caveat for fairness auditing in financial-inclusion modelling more broadly, discussed further in Section~\ref{sec:ethical_considerations}.
 
More substantively, we profiled the feature means of individuals with true label Excluded, comparing those correctly classified against those misclassified as Digitally Engaged (Table~\ref{tab:error_profile}). Misclassified individuals show markedly higher mobile phone ownership (93.5\% vs.\ 49.3\%), higher internet use (41.9\% vs.\ 20.7\%), higher workforce participation (80.6\% vs.\ 28.6\%), and higher income, and are less likely to be female or rural. This profile closely matches the ``Connected but Unengaged'' cluster identified independently in Section~\ref{subsec:h1_access_usage_gap_rq1}, providing a second, model-error-based confirmation of the H1 finding through an entirely different analytical route.
 
\begin{table}[h]
\caption{Feature Means for True-Excluded Cases: Correct vs.\ Misclassified}
\label{tab:error_profile}
\centering
\resizebox{\linewidth}{!}{%
\begin{tabular}{@{}lcc@{}}
\toprule
Feature & Misclassified & Correctly Classified \\
\midrule
\% Female & 12.9\% & 69.3\% \\
Mean age & 38.1 & 33.8 \\
\% In workforce & 80.6\% & 28.6\% \\
\% Rural & 38.7\% & 51.4\% \\
\% Mobile phone (con1) & 93.5\% & 49.3\% \\
\% Internet use & 41.9\% & 20.7\% \\
Mean income quintile & 3.32 & 3.13 \\
\bottomrule
\end{tabular}%
}
\end{table}
 
\subsubsection{Ablation Study}
\label{subsubsec:ablation_study_results}
Table~\ref{tab:ablation} and Fig.~\ref{fig:ablation} show macro F1 under each feature-group configuration. Removing the behaviour group causes the largest performance drop (0.494 $\rightarrow$ 0.415, $-16.1\%$), while removing the access group slightly \emph{improves} macro F1 (0.494 $\rightarrow$ 0.506). Access-only as the sole feature group performs worst of all configurations (macro F1 $=0.186$, barely above a majority-class baseline). This provides a second, methodologically independent confirmation of the SHAP-based ranking in Section~\ref{subsec:h3_predictors_of_adoption_rq3}: behavioural variables carry more predictive weight than raw access variables in this population.
 
\begin{table}[h]
\caption{Ablation Study Results}
\label{tab:ablation}
\centering
\begin{tabular}{@{}lcc@{}}
\toprule
Configuration & $n$ features & Macro F1 \\
\midrule
Full feature set & 15 & 0.494 \\
All except demographics & 9 & 0.480 \\
All except access & 13 & 0.506 \\
All except behaviour & 10 & 0.415 \\
All except resilience & 13 & 0.489 \\
Demographics only & 6 & 0.351 \\
Access only & 2 & 0.186 \\
Behaviour only & 5 & 0.389 \\
Resilience only & 2 & 0.305 \\
\bottomrule
\end{tabular}
\end{table}
 
\begin{figure}[h]
\centering
% TODO: add ablation_study.png here
\includegraphics[width=\linewidth]{ablation_study.png}
\caption{Ablation study: macro F1 under feature-group leave-one-out and single-group-only configurations.}
\label{fig:ablation}
\end{figure}

\subsection{Model Robustness}
\label{subsec:model_robustness_results}
Repeated stratified cross-validation (50 folds) of the tuned Random Forest yielded a macro F1 distribution with mean $=0.480$, std $=0.023$, and 95\% percentile interval $[0.447, 0.521]$ (Fig.~\ref{fig:robustness}), tightly bracketing the single-split test estimate of 0.494 and indicating the reported performance is not an artifact of one particular train-test partition.
 
\begin{figure}[h]
\centering
% TODO: add robustness_f1_distribution.png here
\includegraphics[width=0.82\linewidth]{robustness_f1_distribution.png}
\caption{Distribution of macro F1 across 50 repeated stratified cross-validation folds.}
\label{fig:robustness}
\end{figure}
 
\subsection{Cluster Profiles and Gender Association}
\label{subsec:cluster_profiles_gender_association}
This gender effect noted in Section~\ref{subsubsec:error_analysis_results} is confirmed with substantially larger magnitude at the cluster level. Cluster membership was very strongly associated with gender ($\chi^2 = 650.549$, $p < 0.0001$, $\mathit{df} = 2$, Cramer's V $= 0.807$), and Adoption Tier itself showed a large association with gender ($\chi^2 = 162.080$, $p < 0.0001$, $\mathit{df} = 3$, Cramer's V $= 0.403$) -- both well above the effect sizes observed for urban/rural status. The most digitally excluded cluster was 96.3\% female (95\% CI: 94.1\%--97.7\%), compared to 16.3\% and 10.3\% female in the two more digitally engaged clusters (Table~\ref{tab:clusters}).
 
\begin{table}[h]
\caption{Cluster Profiles (k-means, k=3)}
\label{tab:clusters}
\centering
\resizebox{\linewidth}{!}{%
\begin{tabular}{@{}lccccc@{}}
\toprule
Cluster & $n$ (\%) & \% Female & \% Rural & \% Mobile & \% Excluded \\
\midrule
0 -- Digitally Excluded Women & 431 (43.1\%) & 96.3\% & 51.0\% & 28.3\% & 89.3\% \\
1 -- Connected but Unengaged & 443 (44.3\%) & 16.3\% & 46.0\% & 98.0\% & 66.8\% \\
2 -- Advanced Adopters & 126 (12.6\%) & 10.3\% & 28.6\% & 98.4\% & 0.8\% \\
\bottomrule
\end{tabular}%
}
\end{table}
 
\begin{figure}[h]
\centering
% TODO: add kmeans_k_selection.png here
\includegraphics[width=\linewidth]{kmeans_k_selection.png}
\caption{Silhouette, Davies-Bouldin, and Calinski-Harabasz scores for $k=2$ to $5$.}
\label{fig:k_selection}
\end{figure}

\begin{figure}[h]
\centering
% TODO: add cluster_pca_scatter.png here
\includegraphics[width=\linewidth]{cluster_pca_scatter.png}
\caption{PCA projection of k-means clusters ($k=3$).}
\label{fig:pca_scatter}
\end{figure}

\begin{figure}[h]
\centering
% TODO: add cluster_profile_heatmap.png here
\includegraphics[width=\linewidth]{cluster_profile_heatmap.png}
\caption{Normalized cluster profile heatmap across demographic and access variables.}
\label{fig:cluster_heatmap}
\end{figure}

\subsection{Cluster Validation}
\label{subsec:cluster_validation_results}
Bootstrap resampling (100 replicates) showed the $k=3$ solution is highly stable: mean ARI $=0.904$ (std $=0.041$), 95\% percentile interval $[0.847, 0.982]$ (Fig.~\ref{fig:cluster_stability}). An algorithm-independent cross-check against Ward agglomerative clustering at the same $k$ showed substantial agreement (ARI $=0.740$), indicating the three-cluster structure is not an artifact of the k-means algorithm specifically.
 
\begin{figure}[h]
\centering
% TODO: add cluster_stability_bootstrap.png here
\includegraphics[width=\linewidth]{cluster_stability_bootstrap.png}
\caption{Distribution of Adjusted Rand Index across 100 bootstrap resamples, evaluating stability of the k=3 clustering solution.}
\label{fig:cluster_stability}
\end{figure}
 
Despite the marginal silhouette difference between $k=2$ (0.1675) and
$k=3$ (0.1675), the two solutions are substantively different (ARI $=
0.623$). The $k=2$ solution separates the sample almost entirely along the
``Digitally Excluded Women'' cluster versus everyone else (515 of 516
members of one $k=2$ cluster correspond to the original clusters 1 and 2
combined). What $k=3$ adds is a further split of that ``everyone else''
group into the access-rich-but-unengaged cluster (66.8\% still Excluded)
versus the advanced-adopter cluster (0.8\% Excluded) --- a distinction
directly relevant to the access-usage gap central to H1, and the primary
substantive justification for reporting $k=3$ despite its narrow
quantitative margin over $k=2$.

We are explicit about what this evidence does and does not support. The
near-identical silhouette scores at $k=2$ and $k=3$ indicate that three
clusters are not a dramatically better fit to the data's geometry than
two; silhouette score alone would not strongly favor either solution.
\textbf{We therefore interpret $k=3$ as an analytically useful
segmentation --- one that recovers a policy-relevant distinction the
$k=2$ solution collapses --- rather than as evidence of three sharply
separated natural populations.} The bootstrap stability (ARI $= 0.904$)
and the Ward agglomerative cross-check (ARI $= 0.740$) demonstrate that
the $k=3$ partition is a reproducible way of dividing this sample, not
that the underlying population contains three discrete, well-separated
subgroups in the geometric sense that silhouette score measures.
 
\section{Discussion}
\label{sec:discussion}

The central contribution of this analysis is the identification of a
large, access-rich but adoption-poor population segment, discovered
independently through three separate analytical routes: unsupervised
clustering (\S\ref{subsec:h1_access_usage_gap_rq1}), model error analysis (\S\ref{subsubsec:error_analysis_results}), and feature-group
ablation (\S\ref{subsubsec:ablation_study_results}). These three analyses use different data
representations and different statistical logic, yet converge on the same
conclusion, summarized in Table~\ref{tab:convergence}.

\begin{table}[h]
\caption{Three Independent Signatures of the Access-Usage Gap}
\label{tab:convergence}
\centering
\renewcommand{\arraystretch}{1.3}
\begin{tabular}{p{1.7cm}p{2.9cm}p{2.9cm}}
\toprule
\textbf{Method} & \textbf{What it measures} & \textbf{Result} \\
\midrule
Unsupervised clustering (\S\ref{subsec:h1_access_usage_gap_rq1}) &
$k$-means on access/demographic features, blind to Adoption Tier labels &
44.3\% of sample forms a distinct cluster: 98.0\% mobile, 62.3\% internet,
yet 66.8\% formally Excluded \\
\addlinespace
Model error analysis (\S\ref{subsubsec:error_analysis_results}) &
Feature profile of true-Excluded cases the classifier mistakenly predicts
as Digitally Engaged &
Misclassified group shows far higher mobile ownership (93.5\% vs.\ 49.3\%)
and internet use (41.9\% vs.\ 20.7\%) than correctly classified Excluded
cases \\
\addlinespace
Feature-group ablation (\S\ref{subsubsec:ablation_study_results}) &
Macro F1 change when each construct group is removed from the model &
Removing \emph{behaviour} features drops macro F1 by 16.1\%; removing
\emph{access} features \emph{improves} it --- access alone is the weakest
predictor group (F1 = 0.186) \\
\bottomrule
\end{tabular}
\end{table}

At 44.3\% of the sample, this segment is far larger than the 6.6\%
dormant-account rate would suggest, and it reframes the practical question
facing Pakistan's digital finance sector: the primary obstacle to
inclusion is not, for a substantial share of the population, a lack of
phones or internet connectivity, but something downstream of that access
--- plausibly trust, product design, documentation requirements, or
social/institutional barriers not captured directly in this dataset.
\textbf{We are careful to distinguish what this analysis establishes from
what it does not: the data identify \emph{who} is Connected but Unengaged
and confirm, through three independent methods, that this population is
real and substantial --- but nothing in a cross-sectional survey of
access and demographic variables can identify \emph{why} they remain
unengaged.} The candidate mechanisms above are plausible readings of the
literature, not findings of this study. The segmentation therefore
motivates, rather than resolves, a second-stage mixed-methods
investigation into the mechanisms underlying non-adoption within this
segment --- for instance, qualitative interviews or a targeted follow-up
survey instrument capable of distinguishing trust deficits from product-
design friction from documentation barriers, none of which Findex's
existing variables can separate. The convergence of three independent
methods on the same population segment substantially strengthens
confidence in \emph{who} this population is; it says nothing about
\emph{why} they behave this way, and we do not claim otherwise.

\subsection{Practical Implications}
\label{subsec:practical_implications}

The findings above suggest three practical implications for policy and
product design in Pakistan's digital-finance ecosystem.

\textbf{1) Access alone is not enough.} Inclusion strategy should
separate the pathway from \emph{access} to \emph{adoption} to
\emph{active usage}, rather than treating account opening or connectivity
as sufficient endpoints.

\textbf{2) Targeted intervention for the Connected but Unengaged
segment.} The 44.3\% access-rich but adoption-poor population identified
here may be better addressed through financial education, trust-building,
onboarding assistance, and product redesign than through connectivity
expansion alone.

\textbf{3) Gender-sensitive inclusion design.} The strong gender
association in both cluster membership and Adoption Tier suggests that
generic population-level interventions may miss the most excluded group
and should be complemented with gender-responsive targeting and delivery
mechanisms.

\textbf{Interpretive boundary.} These implications are hypotheses for
intervention, not causal conclusions. They follow from robust
associational patterns in this sample and should be tested prospectively
through intervention studies or quasi-experimental designs before being
treated as policy-certainty claims.

We emphasize that the supervised model's role in this paper is
interrogative rather than predictive. A macro F1 of 0.494 (95\% CI:
[0.447, 0.521] under repeated cross-validation) does not support
individual-level deployment, and we do not recommend this classifier, or
similarly modest classifiers trained on comparably small survey samples,
for operational use in targeting or eligibility decisions (see also
Section~\ref{subsec:risk_of_misuse}). Its value in this study lies instead in the explainability
and diagnostic machinery built around it --- SHAP rankings, permutation
importance, calibration analysis, and error profiling --- each of which
contributed independent evidence toward the segmentation finding above. We
view this study as closer to exploratory behavioural segmentation research
than to applied predictive modelling, and believe the paper is stronger,
not weaker, for making that distinction explicit.

The gender findings compound this picture. Rather than a marginal
demographic control variable, gender shows one of the largest effect sizes
in the entire analysis (Cramer's $V = 0.807$ for cluster membership),
exceeding even the urban/rural divide by a wide margin, and echoes the
gendered impacts of mobile money access documented by Suri and
Jack~\cite{surijack2016} in a causal, longitudinal Kenyan setting. Our
cross-sectional Pakistani data cannot establish causality in either
direction, but the consistency of a large gender effect across two very
different national contexts and methodologies is notable. This diverges
from Saidy and Hassan's Sub-Saharan African findings, where account
ownership and technology access were the most influential predictors in
their models~\cite{saidy2025}, and suggests that regional context is
associated with markedly different predictive structures in
fintech-adoption models --- access-related features dominate in the
Sub-Saharan African setting they study, while behavioural and demographic
features dominate here. We stress that this is a difference in
\emph{which features are predictive}, not a demonstrated difference in
\emph{what causes} adoption in each region; we offer it as a hypothesis
worth testing with causally-identified designs, not as evidence that
access-centric interventions are ineffective in Pakistan. If this
predictive pattern does reflect a genuine behavioural difference, it would
suggest that access-centric interventions calibrated for Sub-Saharan
Africa warrant re-evaluation before being applied to Pakistan, where the
access-usage gap and the gender gap appear entangled with each other
rather than separate problems.

The H2 finding --- urban adoption exceeding rural adoption --- aligns with
the mainstream financial inclusion literature but diverges from Chowdhury
et al.'s Bangladesh result~\cite{chowdhury2026}. Since Chowdhury et al.
measured financial literacy rather than behavioural adoption, this
divergence may reflect a genuine difference between knowledge and
behaviour as outcome variables, rather than a contradiction between the
two country contexts.

The calibration results carry a practical implication distinct from the
accuracy and F1 results typically reported in this literature: a model
with modest discriminative performance (macro F1 $\approx 0.49$) can still
be poorly calibrated for its most policy-relevant class. Here, the model
systematically underestimates the probability of exclusion, which would
understate the scale of the excluded population if predicted
probabilities, rather than hard class labels, were used directly for
resource targeting. This is a methodological point we believe deserves
more attention in applied fintech-adoption modelling generally, where
headline accuracy is often reported without accompanying calibration
diagnostics.

The feature importance consistency results similarly complicate a simple
``top predictor'' narrative. While Gini and SHAP rankings agree strongly,
permutation importance --- arguably the most directly interpretable of the
three, since it measures actual performance degradation --- agrees only
weakly and assigns negative importance to several features, including
gender and rural status. We interpret this as evidence that Random
Forest's internal (Gini) importance and SHAP's coalition-based
attributions can overstate the standalone predictive necessity of a
feature when that feature is correlated with others already carrying the
relevant signal (e.g., gender's correlation with the cluster structure
itself), a distinction between \emph{attribution} and \emph{necessity}
that we believe is under-discussed in applied SHAP-based fintech
literature. As throughout this paper, we treat all three rankings as
statements about model attribution rather than claims about which
features causally drive adoption (Section~\ref{subsec:explainability}).

\section{Ethical Considerations}
\label{sec:ethical_considerations}
\subsection{Fairness and the Accuracy Paradox}
\label{subsec:fairness_accuracy_paradox}
Section~\ref{subsubsec:error_analysis_results} demonstrated that raw subgroup accuracy (0.852 for women vs.\ 0.689 for men) is a misleading fairness metric in this setting, because gender is strongly entangled with the outcome variable itself at the population level. A na\"ive fairness audit relying on aggregate accuracy would conclude the model treats women more favorably, when in fact women are disproportionately concentrated in the largest, easiest-to-predict class. We recommend that fairness audits of financial-inclusion models report subgroup-level macro F1 or per-class recall rather than aggregate accuracy, and we flag this as a specific methodological risk for future work in this domain, since account ownership and gender are correlated by construction in many emerging-market contexts, not just Pakistan.
 
\subsection{Risk of Misuse}
\label{subsec:risk_of_misuse}
Adoption-tier prediction models of the type developed here are intended to support inclusion-oriented policy and product design -- for example, identifying the access-rich, adoption-poor segment as a target for trust-building or product-design interventions. The same predictive infrastructure could, however, be repurposed for exclusionary rather than inclusionary ends: for instance, using predicted low adoption likelihood as an input to credit-scoring or loan-eligibility decisions, which would risk encoding and amplifying existing gender and rural disparities into automated lending decisions. We explicitly do not recommend this model, or models like it, be used for individual-level credit or eligibility decisions; its intended use is aggregate, population-level policy targeting.

\subsection{Human-impact Risk}
\label{subsec:human_impact_risk}
Because the identified segments correlate strongly with gender and other
socioeconomic characteristics, segment labels should not be treated as
immutable characteristics of individuals or used to deny services. We
recommend treating segment assignment as a provisional, context-dependent
signal for outreach design, not as a basis for exclusion, reduced access,
or negative automated decision-making at the individual level.

Table~\ref{tab:model_can_cannot} summarizes the practical scope of this
study's modelling approach, separating valid analytical uses from
inappropriate individual-level or causal interpretations.

\begin{table}[h]
\caption{What the Model Can and Cannot Do}
\label{tab:model_can_cannot}
\centering
\renewcommand{\arraystretch}{1.2}
\begin{tabular}{p{3.7cm}p{3.7cm}}
\toprule
\textbf{The model can} & \textbf{The model cannot} \\
\midrule
Identify statistical associations & Establish causality \\
Characterize population segments & Diagnose individuals \\
Identify influential model features & Identify causal mechanisms \\
Support hypothesis generation & Make credit decisions \\
Motivate targeted research & Determine who ``deserves'' financial services \\
\bottomrule
\end{tabular}
\end{table}
 
\subsection{Calibration and Deployment Caution}
\label{subsec:calibration_deployment_caution}
Given the poor calibration observed for the Excluded class (Section~\ref{subsubsec:calibration_analysis_results}), we caution against using this model's predicted probabilities directly in any downstream resource-allocation formula without recalibration (e.g., Platt scaling or isotonic regression), since the systematic underconfidence identified here would bias such a formula toward underestimating exclusion.
 
\subsection{Data Privacy}
\label{subsec:data_privacy}
All analysis was conducted on de-identified, publicly available World Bank microdata collected under the Global Findex survey's standard informed-consent protocols. No attempt was made, or would be possible, to re-identify individual respondents from the released microdata.
 
\section{Limitations}
\label{sec:limitations}
\begin{itemize}
\item \textbf{Cross-sectional design.} All findings in this paper are
associational. Findex is a single-wave cross-sectional survey, and no
causal identification strategy --- panel data, instrumental variables,
or experimental variation --- is employed here. No causal claims are
made regarding the top predictors of modelled fintech adoption or the mechanisms 
underlying the gender and access-usage gaps identified; see Section~III-G for the
predictor/driver terminology convention we follow throughout to keep
this distinction explicit at the sentence level, not just in this
Limitations section.
\item The Basic tier ($n=21$) is too small for reliable classification by
any model tested; this is a sample-size constraint rather than a
modelling failure and should not be over-interpreted. More broadly, with
$n=1{,}000$ split across four imbalanced classes, \textbf{this model
should not be interpreted as a reliable individual-level classifier} ---
consistent with the modelling objective stated in Section~III-D, its
role in this paper is to surface aggregate, population-level correlates
of adoption through explainability tooling, not to predict any single
respondent's Adoption Tier with confidence. Per-respondent predictions
from this model, including the specific misclassified cases discussed in
Section~IV-C-4, should be read as illustrative of aggregate error
patterns, not as reliable statements about those individuals.
\item The variable \texttt{fin17b} (saved via mobile money), the top-ranked predictor by both Gini and SHAP, is partly confounded with mobile money account ownership due to Findex's survey routing logic (approximately 74\% of respondents were routed out of this question); its top ranking should be interpreted with this caveat.
\item Permutation importance disagreed substantively with Gini and SHAP for several features, including assigning negative importance to gender and rural status; we interpret this as reflecting correlated-feature attribution effects rather than concluding these features are unimportant, but this disagreement itself is a limitation of relying on any single importance method.
\item The three-cluster solution is interpretively and substantively
justified (Section~\ref{subsec:cluster_validation_results}) but is not a dramatically better geometric fit
than $k=2$ by silhouette score alone; see Section~\ref{subsec:cluster_validation_results} for the explicit
distinction we draw between analytical usefulness and natural population
separation.
\item Calibration is poor for the Excluded, Basic, and Digitally Engaged classes; only the Advanced class is well-calibrated. Predicted probabilities for the other three classes should not be used without recalibration.
\item Several chi-square tests involve subgroups with small expected cell counts (notably comparisons involving the Basic tier); results for these specific comparisons should be treated as indicative rather than definitive.
\item Subgroup accuracy is confounded by the strong association between gender and outcome in this sample, as discussed in Section~\ref{subsec:fairness_accuracy_paradox}; this is a general risk for fairness auditing in similar settings, not specific to a modelling error here.
\item \textbf{Generalizability and survey weighting.} All proportions and
model estimates reported in this paper --- including the 44.3\%
``Connected but Unengaged'' segment size, the 68.2\% Excluded rate, and
all cluster and subgroup percentages --- describe the $n=1{,}000$ Findex
Pakistan sample as analyzed, computed without applying Global Findex's
survey design weights. \textbf{These figures should not be interpreted
as population-representative prevalence estimates for Pakistan as a
whole} unless independently reweighted using the Findex sampling weights
and design variables, which we did not do here. Findex microdata
typically includes weights correcting for unequal selection probability
and post-stratification to national demographic totals; because our
pipeline treats each respondent as equally weighted, subgroups
oversampled or undersampled relative to the true population (e.g., by
region, urbanicity, or gender) could shift the reported proportions if
corrected. We recommend that any policy application of these findings
first replicate the segmentation using appropriately weighted estimates.
\end{itemize}
 
\section{Future Work}
\label{sec:future_work}
Several extensions would strengthen this analysis. 
First, replicating the pipeline on the Financial Inclusion Insights (FII) Bangladesh tracker survey, if obtained, would provide a genuine same-instrument cross-country comparison. 
Second, and following directly from the population/mechanism distinction
drawn in Section~V, a qualitative or mixed-methods follow-up focused
specifically on the ``Connected but Unengaged'' cluster could clarify
which downstream barriers --- trust, documentation, product design ---
best explain non-adoption despite access.
Third, model recalibration (Platt scaling or isotonic regression) applied specifically to the Excluded class would address the systematic underconfidence identified in Section~\ref{subsubsec:calibration_analysis_results} and should be evaluated before any applied use of predicted probabilities. Fourth, a fairness-aware modelling approach (e.g., equalized-odds post-processing) could be evaluated as a remedy for the subgroup accuracy confound identified in Section~\ref{subsec:fairness_accuracy_paradox}, alongside subgroup-disaggregated macro F1 reporting as a standard diagnostic. Fifth, causal identification strategies (e.g., difference-in-differences around a policy change, or an instrumental-variables approach) would be needed to move beyond the associational claims made throughout this study.
 
\section{Conclusion}
\label{sec:conclusion}
This study suggests that the next frontier of digital financial inclusion 
in Pakistan may not be connecting people to digital infrastructure, but 
understanding why access fails to translate into participation. The 
identification of a large Connected but Unengaged segment provides a 
measurable target for future mixed-methods and intervention research, 
while the methodological findings demonstrate the importance of leakage 
control, calibration analysis, robustness testing, and careful fairness 
interpretation in financial-inclusion machine learning. Taken together, 
these findings suggest that the next frontier of financial-inclusion research in Pakistan is not simply expanding digital access, but understanding and addressing the gap between having the means to participate and actually participating.
 
\begin{thebibliography}{00}
\bibitem{chowdhury2026} T. A. Chowdhury, M. A. H. Chowdhury, M. T. Rahman, I. Ahmed, N. Ahmed, M. A. I. Tuhin, and A. A. Kafy, ``Modeling financial literacy through explainable machine learning and behavioral segmentation in emerging economies,'' \emph{Computers in Human Behavior Reports}, 2026.
 
\bibitem{emuron2025} A. Emuron, A. D. Kapim Kenfack, and M. Msimanga, ``Fintech predictive modelling using machine learning: An empirical investigation of South Africa,'' \emph{African Journal of Science, Technology, Innovation and Development}, pp.\ 889--899, 2025. doi: 10.1080/20421338.2025.2552006.
 
\bibitem{saidy2025} S. Saidy and A. D. Hassan, ``Modeling drivers of fintech adoption in Sub-Saharan Africa using machine learning,'' \emph{Digital Finance}, vol.\ 7, pp.\ 1065--1091, 2025. doi: 10.1007/s42521-025-00161-4.
 
\bibitem{worldbank2025} World Bank, ``Global Findex Database 2025, Pakistan,'' Reference: PAK\_2024\_FINDEX\_v02\_M, World Bank Microdata Library, 2025. [Online]. Available: \url{https://microdata.worldbank.org/index.php/catalog/7961}
 
\bibitem{surijack2016} T. Suri and W. Jack, ``The long-run poverty and gender impacts of mobile money,'' \emph{Science}, vol.\ 354, no.\ 6317, pp.\ 1288--1292, 2016. doi: 10.1126/science.aah5309.
 
\bibitem{demirguckunt2022} A. Demirg\"u\c{c}-Kunt, L. Klapper, D. Singer, and S. Ansar, \emph{The Global Findex Database 2021: Financial Inclusion, Digital Payments, and Resilience in the Age of COVID-19}. Washington, DC: World Bank, 2022.
 
\end{thebibliography}
 
\appendices
 
\section{Hyperparameter Search Space}
\label{sec:appendix_hyperparameter_search_space}

\begin{table}[h]
\caption{Hyperparameter Search Space and Selected Values}
\label{tab:hyperparam_space}
\centering
\resizebox{\linewidth}{!}{%
\begin{tabular}{@{}lll@{}}
\toprule
Model & Parameter & Search space (selected) \\
\midrule
\multirow{4}{*}{Random Forest} & n\_estimators & \{200, 300, 500, 800\} (300) \\
 & max\_depth & \{4, 6, 8, 10, None\} (None) \\
 & min\_samples\_leaf & \{1, 2, 4, 8\} (4) \\
 & max\_features & \{sqrt, log2, None\} (None) \\
\midrule
\multirow{5}{*}{XGBoost} & n\_estimators & \{200, 300, 500\} (500) \\
 & max\_depth & \{3, 4, 5, 6\} (6) \\
 & learning\_rate & \{0.01, 0.03, 0.05, 0.1\} (0.1) \\
 & subsample & \{0.7, 0.85, 1.0\} (0.85) \\
 & colsample\_bytree & \{0.7, 0.85, 1.0\} (0.85) \\
\bottomrule
\end{tabular}%
}
\end{table}
Table~\ref{tab:hyperparam_space} lists the search space used for \texttt{RandomizedSearchCV} (25 iterations, 5-fold stratified CV, scored on macro F1) and the selected best configuration for each model.
 
Best cross-validated macro F1: Random Forest $=0.472$; XGBoost $=0.457$. Held-out test macro F1: default Random Forest $=0.486$; tuned Random Forest $=0.494$.
 
\section{Additional Diagnostic Tables}
\label{sec:appendix_additional_diagnostic_tables}
Table~\ref{tab:full_importance} reports the complete feature importance comparison across all fifteen features and three methods (Gini, permutation, SHAP), supplementing the abbreviated version in Table~\ref{tab:importance_consistency}.
 
\begin{table}[h]
\caption{Full Feature Importance Comparison}
\label{tab:full_importance}
\centering
\resizebox{\linewidth}{!}{%
\begin{tabular}{@{}lccc@{}}
\toprule
Feature & Gini & Permutation & SHAP \\
\midrule
fin17b & 0.2099 & 0.1087 & 0.0915 \\
age & 0.2307 & 0.0059 & 0.0571 \\
in\_workforce & 0.0786 & 0.0221 & 0.0551 \\
fin17c & 0.0689 & 0.0459 & 0.0364 \\
inc\_q & 0.0905 & 0.0067 & 0.0319 \\
con1 & 0.0437 & $-0.0108$ & 0.0242 \\
fin24b & 0.0808 & 0.0105 & 0.0211 \\
female & 0.0306 & $-0.0169$ & 0.0191 \\
internet\_use & 0.0349 & $-0.0077$ & 0.0187 \\
fin24a & 0.0486 & $-0.0031$ & 0.0180 \\
borrowed & 0.0329 & 0.0125 & 0.0139 \\
educ & 0.0225 & $-0.0038$ & 0.0094 \\
is\_rural & 0.0266 & $-0.0085$ & 0.0077 \\
fin22a\_1 & 0.0005 & $-0.0001$ & 0.0003 \\
merchantpay\_dig & 0.0004 & $-0.0005$ & 0.0003 \\
\bottomrule
\end{tabular}%
}
\end{table}
 
Spearman rank correlations: Gini vs.\ SHAP $\rho = 0.914$ ($p < 0.0001$); Permutation vs.\ SHAP $\rho = 0.504$ ($p = 0.0557$); Gini vs.\ Permutation $\rho = 0.554$ ($p = 0.0323$).
 
Table~\ref{tab:full_confusion} reports the complete confusion matrix for the tuned Random Forest on the held-out test set ($n=250$).
 
\begin{table}[h]
\caption{Full Confusion Matrix (Tuned Random Forest, Test Set)}
\label{tab:full_confusion}
\centering
\begin{tabular}{@{}lcccc@{}}
\toprule
True $\backslash$ Predicted & Excluded & Basic & Dig.\ Engaged & Advanced \\
\midrule
Excluded & 140 & 3 & 21 & 7 \\
Basic & 3 & 0 & 2 & 0 \\
Digitally Engaged & 16 & 0 & 8 & 1 \\
Advanced & 4 & 0 & 2 & 43 \\
\bottomrule
\end{tabular}
\end{table}
 
Table~\ref{tab:cluster_validation_summary} summarizes the cluster validation metrics reported in Section~\ref{subsec:cluster_validation_results}.
 
\begin{table}[h]
\caption{Cluster Validation Summary}
\label{tab:cluster_validation_summary}
\centering
\begin{tabular}{@{}lc@{}}
\toprule
Metric & Value \\
\midrule
Bootstrap ARI, mean (std) & 0.904 (0.041) \\
Bootstrap ARI, 95\% interval & [0.847, 0.982] \\
ARI vs.\ Ward agglomerative & 0.740 \\
ARI, $k{=}2$ vs.\ $k{=}3$ solution & 0.623 \\
Silhouette score, $k{=}2$ & 0.1675 \\
Silhouette score, $k{=}3$ & 0.1675 \\
\bottomrule
\end{tabular}
\end{table}

\vspace{12pt}

\end{document}
